% =====================================================================
% DRAFT FRAGMENT: new Section 3 "Situating the account"
% Built 2026-06-24 for the metaphor-studies retarget (post M&S desk-reject).
% Rev. 3 (2026-06-24 pm): PROJECTIBILITY-FIRST framing (Brett's steer).
%   Lead criterion = projectibility of cross-domain inference (Goodman 1955);
%   mechanism support is the UNDERWRITER, not the headline. The 2x2 axis is
%   relabelled mechanism-support -> projectibility. Goodman glossed + cited at
%   first mention (house style: projectibility needs a gloss for non-phil-sci
%   readers). Goodman1955 is in the CENTRAL bib already.
% Primary-grounded (rev.2): Kövecses 2017 / 2020-article / Murphy & Medin 1985,
%   all page-verified via pdftotext. Steen via Steen 2008 + Flusberg 2018.
% NOT yet in main.tex. First-person voice is Brett's to finalise.
% Macro check: \metaphor{} defined in main.tex preamble (uppercases). \term,
%   \mention, \enquote, \citep, \textcite, \citealt per house style.
% =====================================================================

\section{Situating the account}
\label{sec:situating}

Conceptual Metaphor Theory has not stood still since 1980, and two developments speak directly to the gap I've described. Neither closes it, but saying why locates the contribution. The question this account puts at the centre is the projectibility of metaphorical inference: what observing some features of a mapping licenses us to predict about the rest \citep{Goodman1955}. I take \term{projectibility} in the profile sense of \textcite{reynolds2026kindsProjectibilityProfiles}: a category earns standing by tracking a worldly pattern whose partial observation licenses inferences under specified conditions. Which cross-domain mappings support reliable inference, and what underwrites that reliability? Extended conceptual metaphor theory and deliberate metaphor theory answer different questions, and the three can be held apart.

\subsection{Extended conceptual metaphor theory}

\textcite{kovecses2017, kovecses2020} extends CMT along two axes. First, a single conceptual metaphor is reconstrued as a hierarchy across four levels of schematicity~-- image schema, domain, frame, and mental space~-- rather than a mapping between two domains \citep{kovecses2017}. This dissolves a long-standing grain problem: source domains carry more material than any mapping uses, and locating the metaphor on the hierarchy says which material is in play. Second, metaphor in discourse is given a contextual component. Kövecses distinguishes four context types (situational, discourse, bodily, and conceptual-cognitive), each with factors that \enquote{prime} a speaker toward one metaphor rather than another in a given situation \citep{kovecses2015, kovecses2020}. This answers a real complaint: classical CMT, built from dictionaries and decontextualized examples, had little to say about why a particular metaphor surfaces in particular discourse.

Both moves operate within the synchronic life of a metaphor, and neither asks whether the mapping supports projectible inference. The contextual component is a theory of production-side selection: which stored mapping a speaker reaches for, here and now, under contextual pressure. And the dynamism Kövecses adds is the dynamism of processing, not of history. The offline/online distinction in extended CMT is a distinction between memory systems: image schemas, domains, and frames are \enquote{conventionalized knowledge structures in long-term memory,} while mental spaces are \enquote{online representations of our understanding of experience in working memory} \citep[326]{kovecses2017}. Conventionalization enters as a property a structure already has (stored, shared across a speech community), not as something the theory explains, and reliability of inference doesn't enter at all. Extended CMT takes the supraindividual layer~-- the decontextualized patterns a language and culture happen to share~-- as given, and theorizes what speakers do with it online. Which of those mappings project, and what keeps them projectible across generations, is a different question. Mechanism support answers it: a mapping projects when the shared relational structure is causally maintained in both domains, not just in the source. The two layers nest rather than compete. Contextual priming selects among mappings that are already available; projectibility, underwritten by mechanisms, is what made them worth keeping.

The difference is testable, not terminological. Extended CMT doesn't predict that peripheral, weakly maintained properties erode before core ones when transmission weakens (the Dyirbal and \mention{data} cases, Section~\ref{sec:source-domains}); it doesn't predict that a high-frequency mapping fails to conventionalize when its target domain has no mechanisms sustaining the shared structure (Prediction~1); and it makes no prediction about non-linguistic animals (Section~\ref{sec:dogs}). The projectibility account does. And the gap is one Kövecses partly concedes from the inside: his processing model is, in his own words, \enquote{informal and hypothetical,} and \enquote{it is only through psychological experiments that the validity of this \ldots\ model could be confirmed} \citep[125]{kovecses2020article}. Grounding reliability of inference in world-side mechanisms is one way to give the framework external traction it doesn't yet have.

\subsection{Deliberate metaphor}

\textcite{steen2008} adds a third dimension to the linguistic and conceptual ones: a metaphor can be used \emph{deliberately}, as a metaphor that points the addressee to a source domain, or non-deliberately, as part of the ordinary structure of the language, where \metaphor{argument is war} passes unnoticed. Deliberateness is a fact about how a metaphor functions in communication. It cross-cuts projectibility rather than tracking it, and the two axes are worth separating explicitly.

\begin{table}[ht]
\centering\small
\begin{tabular}{@{}>{\raggedright\arraybackslash}p{.20\textwidth}>{\raggedright\arraybackslash}p{.36\textwidth}>{\raggedright\arraybackslash}p{.36\textwidth}@{}}
\toprule
 & Deliberate & Non-deliberate \\
\midrule
High projectibility & A theory-constitutive metaphor a scientist knowingly deploys (\metaphor{the brain is a computer}) to direct research & A constitutive mapping so entrenched it's gone invisible in the field that lives by it \\[6pt]
Low projectibility & A fresh literary metaphor (Dickinson's \enquote{Hope is the thing with feathers}) & A dead idiom whose source no longer projects anything \\
\bottomrule
\end{tabular}
\caption{Deliberateness (Steen) and projectibility cross-cut. Mechanism support underwrites the rows, not the columns.}
\label{tab:deliberate-projectibility}
\end{table}

The three-way typology developed below (literary, everyday, constitutive; Section~\ref{sec:literary}) is the projectibility axis of Table~\ref{tab:deliberate-projectibility}; deliberateness is the other. Because the axes are independent, whatever processing signature deliberateness predicts, projectibility predicts a separate one: how reliably the mapping supports inference, how gracefully it degrades when source-domain knowledge is disrupted, and whether it conventionalizes over time. The two can be measured apart. That extended CMT itself lists incorporating deliberate metaphor among its open problems \citep[113]{kovecses2020article}, rather than deriving it from its apparatus, is a sign the communicative axis really is separate from the conceptual one. I take the distinction DMT draws, not its further claim that deliberate metaphors capture more attention and are better remembered~-- an entailment that \textcite{gibbs2015dmt} and \textcite{thibodeau2017} find poorly supported, as \textcite{flusberg2018} note.

\subsection{Cultural and cross-linguistic variation}

Communities with the same bodies metaphorize differently, and the differences are patterned, not free (Section~\ref{sec:source-domains}). Embodiment alone struggles here: if shared sensorimotor experience grounds the mappings, shared bodies should converge. The projectibility account locates the variation in the mechanisms rather than the bodies. Where ecologies, institutions, and practices differ across communities, the property clusters they maintain differ too, and so do the mappings that stay projectible. This gives a mechanism reading of the cultural variation \textcite{kovecses2005} documents: variation tracks differences in what the world, locally, keeps clustered, which also predicts which parts of a mapping should project across communities and which should fray. The point connects to discourse-based and cross-linguistic metaphor analysis \citep[e.g.][]{semino2008, musolff2006}, which shows metaphors doing situated cultural work that a purely embodied account underdetermines. Projectibility says which of those situated mappings have the causal backing to support reliable inference and persist.

\subsection{Concept structure and analogy}

The account already treats structure-mapping systematicity \citep{gentner1983} and the purpose constraint \citep{holyoak2018} as components of cross-domain extension (Section~\ref{sec:consequences}). It connects to an older claim about category structure. \textcite{murphymedin1985} argued that what makes a concept cohere isn't similarity among its features but the background theories a learner brings: \enquote{concepts are coherent to the extent that they fit people's background knowledge or naive theories about the world,} theories that \enquote{help to relate the concepts in a domain and to structure the attributes that are internal to a concept} \citep[289]{murphymedin1985}. Similarity, they argue, is too unconstrained to pick out which concepts are coherent~-- which is, in Goodman's terms, the projectibility problem for concepts. HPC is the world-side completion of that move. A learner's theory projects to new instances to the degree it tracks causal mechanisms in the world that maintain the cluster. This separates the account from similarity-based models (prototype and exemplar), where cohesion lives in the representation, and from structure-mapping read as purely formal alignment: what makes an alignment support reliable inference is causal maintenance in both domains, not the match itself.
